\section{Winkel -- Gradmaß und Bogenmaß}
\label{sec:fundamentals_angle_measures}

Winkel beschreiben Drehungen in der Ebene. Je nachdem, ob eine vollständige
Drehung mit $360^\circ$ oder mit $2\pi$ angegeben wird, verwenden wir das
Gradmaß oder das Bogenmaß.

\begin{DefinitionBox}{Winkel}{fundamentals_angle}
  Ein \emph{Winkel} \EnglishTerm{angle} beschreibt die Größe einer Drehung
  eines Strahls um seinen Anfangspunkt. Der Winkel zwischen der Ausgangs- und
  der Endlage wird üblicherweise mit einem griechischen Buchstaben wie
  $\alpha$ bezeichnet.

  \begin{center}
    \begin{tikzpicture}[scale=1.05,line cap=round,line join=round]
  \coordinate (O) at (0,0);
  \draw[very thick,->] (O) -- (4.0,1.75);
  \draw[very thick,->] (O) -- (4.2,-1.25);
  \draw[THMBlue,very thick] (1.15,-0.34) arc[start angle=-16.6,end angle=23.6,radius=1.21];
  \node[THMBlue,font=\large] at (0.72,0.04) {$\alpha$};
  \fill (O) circle (1.5pt);
\end{tikzpicture}

  \end{center}
\end{DefinitionBox}

\begin{DefinitionBox}{Gradmaß und Bogenmaß}{fundamentals_degree_and_radian_measure}
  Beim \emph{Gradmaß} \EnglishTerm{degree measure} wird eine vollständige
  Drehung in $360$ gleich große Teile unterteilt. Ein voller Winkel besitzt
  daher $360^\circ$, ein rechter Winkel $90^\circ$.

  Beim \emph{Bogenmaß} \EnglishTerm{radian measure} entspricht eine
  vollständige Drehung dem Wert $2\pi$. Ein rechter Winkel besitzt damit den
  Wert $\frac{\pi}{2}$.
\end{DefinitionBox}

Die folgende Gegenüberstellung zeigt wichtige Winkel in beiden Winkelmaßen:

\begin{center}
  \begin{tikzpicture}[
    x=2.25cm,
    y=1.0cm,
    line cap=round,
    every node/.style={font=\small}
  ]
  \draw[->,thick] (-0.35,0.55) -- (4.55,0.55)
    node[right] {$\alpha_{\mathrm G}$};
  \draw[->,thick,THMBlue] (-0.35,-0.75) -- (4.55,-0.75)
    node[right] {$\alpha_{\mathrm B}$};

  \foreach \x/\degree/\radian in {
      0/{0^\circ}/{0},
      1/{90^\circ}/{\frac{\pi}{2}},
      2/{180^\circ}/{\pi},
      3/{270^\circ}/{\frac{3\pi}{2}},
      4/{360^\circ}/{2\pi}
    } {
      \draw[black!45,densely dashed] (\x,1.12) -- (\x,-1.18);
      \draw[thick] (\x,0.45) -- (\x,0.65);
      \draw[THMBlue,thick] (\x,-0.85) -- (\x,-0.65);
      \node[below] at (\x,0.42) {$\degree$};
      \node[below,THMBlue] at (\x,-0.88) {$\radian$};
    }

  \node[anchor=east,font=\bfseries] at (-0.45,0.55)
    {\Lang{Gradmaß}{Degree measure}:};
  \node[anchor=east,font=\bfseries,THMBlue] at (-0.45,-0.75)
    {\Lang{Bogenmaß}{Radian measure}:};

  \node[above,text=black!60,align=center] at (0,1.12)
    {\Lang{keine\\Drehung}{no\\rotation}};
  \node[above,text=black!60,align=center] at (1,1.12)
    {\Lang{Viertel-\\drehung}{quarter\\turn}};
  \node[above,text=black!60,align=center] at (2,1.12)
    {\Lang{halbe\\Drehung}{half\\turn}};
  \node[above,text=black!60,align=center] at (3,1.12)
    {\Lang{Dreiviertel-\\drehung}{three-quarter\\turn}};
  \node[above,text=black!60,align=center] at (4,1.12)
    {\Lang{ganze\\Drehung}{full\\turn}};
\end{tikzpicture}

\end{center}

\begin{KeyPointBox}{Umrechnung zwischen Gradmaß und Bogenmaß}{fundamentals_degree_radian_conversion}
  Da $360^\circ$ genau $2\pi$ im Bogenmaß entsprechen, gilt
  \begin{ClozeArea}[0.82\linewidth]{3cm}
    \begin{align*}
      \frac{\alpha_{\mathrm G}}{360^\circ}
        &=\frac{\alpha_{\mathrm B}}{2\pi},\\
      \alpha_{\mathrm B}
        &=\alpha_{\mathrm G}\cdot\frac{\pi}{180^\circ},
      &
      \alpha_{\mathrm G}
        &=\alpha_{\mathrm B}\cdot\frac{180^\circ}{\pi}.
    \end{align*}
  \end{ClozeArea}
\end{KeyPointBox}

\begin{ExampleBox}{Umrechnung von Winkeln}{fundamentals_angle_conversion_example}
  \begin{enumerate}[label=\alph*)]
    \item $60^\circ$ entsprechen im Bogenmaß
      \[
        60^\circ\cdot\frac{\pi}{180^\circ}
          =\Cloze{2.5cm}{\frac{\pi}{3}}.
      \]
    \item $\frac{3\pi}{4}$ entsprechen im Gradmaß
      \[
        \frac{3\pi}{4}\cdot\frac{180^\circ}{\pi}
          =\Cloze{2.5cm}{135^\circ}.
      \]
  \end{enumerate}
\end{ExampleBox}

Das Bogenmaß besitzt eine direkte geometrische Bedeutung. Ein Winkel
$\alpha_{\mathrm B}$ schneidet aus einem Kreis mit Radius $R$ einen Kreisbogen
der Länge $b$ aus.

\begin{center}
  \begin{tikzpicture}[scale=1.25,line cap=round,line join=round]
  \def\radius{2.0}
  \def\angle{52}

  \fill[THMBlue!8] (0,0) -- (\radius,0)
    arc[start angle=0,end angle=\angle,radius=\radius] -- cycle;
  \draw[black!45,thick] (0,0) circle (\radius);
  \draw[very thick] (0,0) -- (\radius,0);
  \draw[very thick] (0,0) -- ({\radius*cos(\angle)},{\radius*sin(\angle)});
  \draw[THMBlue,very thick]
    (\radius,0) arc[start angle=0,end angle=\angle,radius=\radius];
  \draw[THMBlue,thick]
    (0.85,0) arc[start angle=0,end angle=\angle,radius=0.85];

  \node[THMBlue] at ({0.48*cos(\angle/2)},{0.48*sin(\angle/2)}) {$\alpha$};
  \node[above left] at ({1.22*cos(\angle)},{1.22*sin(\angle)}) {$R$};
  \node[THMBlue,right] at ({\radius*cos(\angle/2)+0.08},{\radius*sin(\angle/2)}) {$b$};
  \fill (0,0) circle (1.4pt);
\end{tikzpicture}

\end{center}

\begin{KeyPointBox}{Bogenlänge eines Kreisausschnitts}{fundamentals_circular_arc_length}
  Für einen Winkel $\alpha_{\mathrm B}$ im Bogenmaß und einen Kreis mit Radius
  $R$ gilt
  \begin{ClozeArea}[0.65\linewidth]{2cm}
    \[
      b=\alpha_{\mathrm B}\,R.
    \]
  \end{ClozeArea}
  Für den Einheitskreis mit $R=1$ ist die Bogenlänge daher gleich dem
  Winkelwert: $b=\alpha_{\mathrm B}$. Für eine vollständige Drehung ergibt sich
  als Spezialfall der Kreisumfang $U=2\pi R$.
\end{KeyPointBox}

\begin{RemarkBox}
  Winkel im Bogenmaß sind einheitenlos; die Schreibweise ``rad'' dient nur der
  Kennzeichnung. Beim Gradmaß gehört dagegen das Gradzeichen $^\circ$ zum
  Winkelwert. In Mathematik und Physik wird daher meist das Bogenmaß verwendet.
\end{RemarkBox}
